% GENERATED FILE. Edit canonical lessons, then run npm run generate:books.
% canonical-source-hash: fnv1a64:f5ebde245583dd47
% canonical-lessons: TE-C459-phalaharam, TE-C459-bhojanasala, TE-C459-sadupayam, TE-C459-antarjalam, TE-C459-charjaru

\chapter{Meals and Comforts}
\label{ch:te-meals-and-comforts}
\hlchaptermodality{459}

\begin{chapteropening}
\textbf{By the end of this chapter:} \emph{I can say a snack, a dining hall, a facility, the web and a charger.}

Complete the last lesson of chapter 459: I can say a snack, a dining hall, a facility, the web and a charger.
\end{chapteropening}

\section[\texorpdfstring{\te{ఫలహారం} (phalahāraṁ)}{ఫలహారం (phalahāraṁ)}]{\te{ఫలహారం} (phalahāraṁ) — a snack; a light meal}

\label{lesson:TE-C459-phalaharam}



\begin{quote}
Before the new one: say the Telugu for a toilet, then the Telugu for drinking water.
\end{quote}

\subsection*{You'll want to know: \te{ఫలహారం}}
\textbf{\te{ఫలహారం}} — \emph{phalahāraṁ} — \textquotedblleft{}a snack; a light meal\textquotedblright{}.

\begin{quote}
\te{ప్రయాణంలో} \te{అమ్మ} \te{ఫలహారం} \te{తెచ్చింది}.
\end{quote}

\begin{quote}
\emph{prayāṇaṁlō amma phalahāraṁ teccindi.}
\end{quote}

\begin{quote}
Mother brought a snack for the journey.
\end{quote}

In Sanskrit it meant living on fruit; in Telugu it is any light meal, such as tiffin.

\begin{cousinweb}
\textbf{\te{ఫలహారం}} (\emph{phalahāraṁ}) is borrowed from Sanskrit \emph{phalāhāra}, \textquotedblleft{}feeding on fruit\textquotedblright{} (Monier-Williams), from \emph{phala}, \textquotedblleft{}fruit\textquotedblright{}, and \emph{āhāra}, \textquotedblleft{}food\textquotedblright{}.
\end{cousinweb}

\begin{grammarlens}[title={What you've built}]
One more word for getting around.
\end{grammarlens}

\subsection*{Your turn}
\begin{itemize}
\raggedright
  \item \emph{Say it:} \emph{phalahāraṁ}
  \item \emph{Say it:} \emph{phalahāraṁ}, once more
  \item \emph{Say it:} \emph{prayāṇaṁlō amma phalahāraṁ teccindi.} — Mother brought a snack for the journey.
  \item \emph{Recall:} say the Telugu for a toilet, then the Telugu for drinking water, then say \emph{phalahāraṁ} again
\end{itemize}

\begin{tcolorbox}[breakable,colback=teal!4,colframe=teal!35!black,title={Before you move on}]
What does \te{ఫలహారం} mean? (\textquotedblleft{}A snack; a light meal\textquotedblright{}.) Say it once more.
\end{tcolorbox}
\section[\texorpdfstring{\te{భోజనశాల} (bhōjanaśāla)}{భోజనశాల (bhōjanaśāla)}]{\te{భోజనశాల} (bhōjanaśāla) — a dining hall}

\label{lesson:TE-C459-bhojanasala}



\begin{quote}
Before the new one: say the Telugu for drinking water, then the Telugu for a snack.
\end{quote}

\subsection*{You'll want to know: \te{భోజనశాల}}
\textbf{\te{భోజనశాల}} — \emph{bhōjanaśāla} — \textquotedblleft{}a dining hall\textquotedblright{}.

\begin{quote}
\te{హోటలు} \te{భోజనశాల} \te{కింది} \te{అంతస్తులో} \te{ఉంది}.
\end{quote}

\begin{quote}
\emph{hōṭalu bhōjanaśāla kindi antastulō undi.}
\end{quote}

\begin{quote}
The hotel dining hall is on the ground floor.
\end{quote}

\te{భోజనం}, a meal, + \te{శాల}, a hall.

\begin{cousinweb}
\textbf{\te{భోజనశాల}} (\emph{bhōjanaśāla}) is a compound of two words borrowed from Sanskrit: \emph{bhojana}, \textquotedblleft{}eating, a meal\textquotedblright{}, and \emph{śālā}, \textquotedblleft{}a shed, a hall\textquotedblright{} (CDIAL 12414).
\end{cousinweb}

\begin{grammarlens}[title={What you've built}]
One more word for getting around.
\end{grammarlens}

\subsection*{Your turn}
\begin{itemize}
\raggedright
  \item \emph{Say it:} \emph{bhōjanaśāla}
  \item \emph{Say it:} \emph{bhōjanaśāla}, once more
  \item \emph{Say it:} \emph{hōṭalu bhōjanaśāla kindi antastulō undi.} — The hotel dining hall is on the ground floor.
  \item \emph{Recall:} say the Telugu for drinking water, then the Telugu for a snack, then say \emph{bhōjanaśāla} again
\end{itemize}

\begin{tcolorbox}[breakable,colback=teal!4,colframe=teal!35!black,title={Before you move on}]
What does \te{భోజనశాల} mean? (\textquotedblleft{}A dining hall\textquotedblright{}.) Say it once more.
\end{tcolorbox}
\section[\texorpdfstring{\te{సదుపాయం} (sadupāyaṁ)}{సదుపాయం (sadupāyaṁ)}]{\te{సదుపాయం} (sadupāyaṁ) — a facility; an amenity}

\label{lesson:TE-C459-sadupayam}



\begin{quote}
Before the new one: say the Telugu for a snack, then the Telugu for a dining hall.
\end{quote}

\subsection*{You'll want to know: \te{సదుపాయం}}
\textbf{\te{సదుపాయం}} — \emph{sadupāyaṁ} — \textquotedblleft{}a facility; an amenity\textquotedblright{}.

\begin{quote}
\te{ఆ} \te{హోటలులో} \te{అన్ని} \te{సదుపాయాలు} \te{ఉన్నాయి}.
\end{quote}

\begin{quote}
\emph{ā hōṭalulō anni sadupāyālu unnāyi.}
\end{quote}

\begin{quote}
That hotel has every amenity.
\end{quote}

Plural: \te{సదుపాయాలు}, \emph{sadupāyālu}.

Plural \te{సదుపాయాలు}, the -\te{ం} noun plural. \te{సౌకర్యం} is a near twin.

\begin{cousinweb}
\textbf{\te{సదుపాయం}} (\emph{sadupāyaṁ}) is built from two Sanskrit words: \emph{sat}, \textquotedblleft{}good\textquotedblright{}, and \emph{upāya}, \textquotedblleft{}a means, an approach\textquotedblright{} (CDIAL 2308): a good means to an end.
\end{cousinweb}

\begin{grammarlens}[title={What you've built}]
One more word for getting around.
\end{grammarlens}

\subsection*{Your turn}
\begin{itemize}
\raggedright
  \item \emph{Say it:} \emph{sadupāyaṁ}
  \item \emph{Say it:} \emph{sadupāyaṁ}, once more
  \item \emph{Say it:} \emph{ā hōṭalulō anni sadupāyālu unnāyi.} — That hotel has every amenity.
  \item \emph{Recall:} say the Telugu for a snack, then the Telugu for a dining hall, then say \emph{sadupāyaṁ} again
\end{itemize}

\begin{tcolorbox}[breakable,colback=teal!4,colframe=teal!35!black,title={Before you move on}]
What does \te{సదుపాయం} mean? (\textquotedblleft{}A facility; an amenity\textquotedblright{}.) Say it once more.
\end{tcolorbox}
\section[\texorpdfstring{\te{అంతర్జాలం} (antarjālaṁ)}{అంతర్జాలం (antarjālaṁ)}]{\te{అంతర్జాలం} (antarjālaṁ) — the web; the internet}

\label{lesson:TE-C459-antarjalam}



\begin{quote}
Before the new one: say the Telugu for a dining hall, then the Telugu for a facility.
\end{quote}

\subsection*{You'll want to know: \te{అంతర్జాలం}}
\textbf{\te{అంతర్జాలం}} — \emph{antarjālaṁ} — \textquotedblleft{}the web; the internet\textquotedblright{}.

\begin{quote}
\te{ఆ} \te{సత్రంలో} \te{అంతర్జాలం} \te{లేదు}.
\end{quote}

\begin{quote}
\emph{ā satraṁlō antarjālaṁ lēdu.}
\end{quote}

\begin{quote}
There is no internet at that inn.
\end{quote}

A formal word made for the internet; in speech people mostly say \te{ఇంటర్నెట్}.

\begin{cousinweb}
\textbf{\te{అంతర్జాలం}} (\emph{antarjālaṁ}) is a modern Telugu coinage built from two Sanskrit words: \emph{antar}, \textquotedblleft{}between, inside\textquotedblright{}, and \emph{jāla}, \textquotedblleft{}a net\textquotedblright{} (CDIAL 5213).
\end{cousinweb}

\begin{grammarlens}[title={What you've built}]
One more word for getting around.
\end{grammarlens}

\subsection*{Your turn}
\begin{itemize}
\raggedright
  \item \emph{Say it:} \emph{antarjālaṁ}
  \item \emph{Say it:} \emph{antarjālaṁ}, once more
  \item \emph{Say it:} \emph{ā satraṁlō antarjālaṁ lēdu.} — There is no internet at that inn.
  \item \emph{Recall:} say the Telugu for a dining hall, then the Telugu for a facility, then say \emph{antarjālaṁ} again
\end{itemize}

\begin{tcolorbox}[breakable,colback=teal!4,colframe=teal!35!black,title={Before you move on}]
What does \te{అంతర్జాలం} mean? (\textquotedblleft{}The web; the internet\textquotedblright{}.) Say it once more.
\end{tcolorbox}
\section[\texorpdfstring{\te{ఛార్జరు} (chārjaru)}{ఛార్జరు (chārjaru)}]{\te{ఛార్జరు} (chārjaru) — a phone charger}

\label{lesson:TE-C459-charjaru}



\begin{quote}
Before the new one: say the Telugu for a facility, then the Telugu for the web.
\end{quote}

\subsection*{You'll want to know: \te{ఛార్జరు}}
\textbf{\te{ఛార్జరు}} — \emph{chārjaru} — \textquotedblleft{}a phone charger\textquotedblright{}.

\begin{quote}
\te{నా} \te{ఛార్జరు} \te{ఇంట్లోనే} \te{మర్చిపోయాను}.
\end{quote}

\begin{quote}
\emph{nā chārjaru iṇṭlōnē marcipōyānu.}
\end{quote}

\begin{quote}
I left my charger at home.
\end{quote}

It takes -\te{ఉ} like other English loans.

\begin{cousinweb}
\textbf{\te{ఛార్జరు}} (\emph{chārjaru}) is a loan from English \emph{charger}, with the Telugu final -\te{ఉ}.
\end{cousinweb}

\begin{grammarlens}[title={What you've built}]
One more word for getting around.
\end{grammarlens}

\subsection*{Your turn}
\begin{itemize}
\raggedright
  \item \emph{Say it:} \emph{chārjaru}
  \item \emph{Say it:} \emph{chārjaru}, once more
  \item \emph{Say it:} \emph{nā chārjaru iṇṭlōnē marcipōyānu.} — I left my charger at home.
  \item \emph{Recall:} say the Telugu for a facility, then the Telugu for the web, then say \emph{chārjaru} again
\end{itemize}

\begin{tcolorbox}[breakable,colback=teal!4,colframe=teal!35!black,title={Before you move on}]
What does \te{ఛార్జరు} mean? (\textquotedblleft{}A phone charger\textquotedblright{}.) Say it once more.
\end{tcolorbox}
